\documentclass[]{article}
\usepackage[margin=1in]{geometry}
%opening
\title{Implicit assessment in psychology: Myths, doubts, and practical applications}
\author{Ottavia M. Epifania}
\date{ottavia.epifania@unipd.it \\ marinaottavia.epifania@unicatt.it}

\renewcommand\baselinestretch{1.5}

\begin{document}

\maketitle

\begin{abstract}
Throughout the past 20 years, the implicit assessment of psychological constructs has become vastly popular in social sciences. But what does implicit mean? In this seminar, an excursus on the meaning and implications of the term implicit will be provided, along with an introduction on the most common implicit measures, such as the Implicit Association Test (IAT) and the Go/No-go Association Task (GNAT). 
The scientific evidence and theoretical thinking that lead to the shift from the identification of implicit with the unconscious to its identification with indirect measurement will be given particular attention.
A practical application with the IAT will be carried out. Specifically, the use of software Inquisit for developing  the IAT will be illustrated.  The IAT data  will be collected and analyzed through the \texttt{implicitMeasures} package in \texttt{R}.

\textbf{Keywords:} Implicit assessment; Implicit social cognition; Implicit Association Test; \texttt{R}
\end{abstract}

\section*{Bio}
Ottavia M. Epifania is a post-doc fellow at the University of Padova, where she is currently working on new Item Response Theory (IRT) procedures for the development of test short forms. 
Her main research interests concern IRT and Rasch modeling of accuracy and time response, implicit assessment, and statistical modeling of implicit measures data.
\end{document}
